% =========================================================================
% UNIT 5: SAMPLING DISTRIBUTIONS
% =========================================================================
\chapter{Unit 5: Sampling Distributions \& The Central Limit Theorem}

\section{Unit Overview \& Exam Weight}
Unit 5 accounts for \textbf{7\% to 12\%} of the AP Statistics Exam. It is the crucial conceptual bridge between probability models and formal statistical inference. Understanding the sampling distributions of proportions ($\hat{p}$) and means ($\bar{x}$) is required to justify every confidence interval and significance test in Units 6 through 9.

\subsection*{Key Unit Vocabulary}
\begin{description}[leftmargin=1.5em, style=nextline]
  \item[Parameter:] A numerical value describing a characteristic of an entire population. Fixed, but typically unknown ($\mu, \sigma, p$).
  \item[Statistic:] A numerical value computed from sample data, used to estimate a population parameter ($\bar{x}, s_x, \hat{p}$).
  \item[Sampling Distribution:] The theoretical probability distribution of all possible values of a statistic taken across all possible random samples of the exact same size $n$ from the same population.
  \item[Sampling Variability:] The fundamental reality that different random samples produce different sample statistics.
  \item[Unbiased Estimator:] A statistic whose sampling distribution has a mean equal to the true value of the population parameter being estimated ($E(\hat{p}) = p$ and $E(\bar{x}) = \mu$).
  \item[Central Limit Theorem (CLT):] For a sufficiently large random sample ($n \ge 30$), the sampling distribution of the sample mean $\bar{x}$ is approximately normal, regardless of the underlying shape of the population distribution!
  \item[10\% Condition:] When sampling without replacement, the sample size must not exceed 10\% of the population ($n \le 0.10N$) to satisfy the independence assumption in standard error formulas.
  \item[Large Counts Condition:] For proportions, both $np \ge 10$ and $n(1-p) \ge 10$ must be satisfied to ensure the sampling distribution of $\hat{p}$ is approximately normal.
\end{description}

\section{Core Concept Lessons}

\begin{lessonbox}[Lesson 5.1: The Three Distributions Distinguished]
To avoid disastrous confusion on the AP exam, strictly distinguish between:
\begin{enumerate}[leftmargin=1.5em]
  \item \textbf{Population Distribution:} Gives the values of the variable for all individuals in the entire population. Fixed mean $\mu$ and spread $\sigma$.
  \item \textbf{Sample Distribution:} Gives the values of the variable for the $n$ individuals in a single selected sample. Summary statistics are $\bar{x}$ and $s_x$.
  \item \textbf{Sampling Distribution:} The distribution of values taken by the statistic (e.g., $\bar{x}$ or $\hat{p}$) across \textbf{all possible samples of size $n$}. As $n$ increases, its spread shrinks!
\end{enumerate}
\end{lessonbox}

\begin{lessonbox}[Lesson 5.2: Sampling Distribution of a Sample Proportion ($\hat{p}$)]
Let $\hat{p} = X/n$ be the sample proportion of successes in an SRS of size $n$ from a population with true proportion $p$:
\begin{equation}
  \text{Center: } \mu_{\hat{p}} = p \quad (\hat{p} \text{ is an unbiased estimator of } p)
\end{equation}
\begin{equation}
  \text{Spread: } \sigma_{\hat{p}} = \sqrt{\frac{p(1-p)}{n}} \quad (\text{Requires 10\% condition: } n \le 0.10N)
\end{equation}
\textbf{Shape (Normality Condition):}
The sampling distribution of $\hat{p}$ is approximately normal if and only if the \textbf{Large Counts Condition} is satisfied:
\begin{equation}
  np \ge 10 \quad \text{AND} \quad n(1-p) \ge 10
\end{equation}
\end{lessonbox}

\begin{lessonbox}[Lesson 5.3: Sampling Distribution of a Sample Mean ($\bar{x}$)]
Let $\bar{x}$ be the sample mean of an SRS of size $n$ drawn from a population with mean $\mu$ and standard deviation $\sigma$:
\begin{equation}
  \text{Center: } \mu_{\bar{x}} = \mu \quad (\bar{x} \text{ is an unbiased estimator of } \mu)
\end{equation}
\begin{equation}
  \text{Spread: } \sigma_{\bar{x}} = \frac{\sigma}{\sqrt{n}} \quad (\text{Requires 10\% condition: } n \le 0.10N)
\end{equation}
\textbf{Shape (Normality Condition):}
\begin{itemize}[leftmargin=1.2em]
  \item If the population distribution is \textbf{Normal}, the sampling distribution of $\bar{x}$ is \textbf{Normal for any sample size $n$}.
  \item If the population distribution is \textbf{NOT Normal (skewed or multimodal)}, the \textbf{Central Limit Theorem (CLT)} guarantees that the sampling distribution of $\bar{x}$ is approximately normal provided the sample size is sufficiently large:
  \begin{equation}
    n \ge 30
  \end{equation}
\end{itemize}
\end{lessonbox}

\begin{trapbox}[Exam Trap Alert: Central Limit Theorem vs. Large Counts]
\begin{itemize}[leftmargin=1.2em]
  \item The \textbf{Central Limit Theorem applies ONLY to sample means ($\bar{x}$)}!
  \item Never write \textit{"The sampling distribution of $\hat{p}$ is normal by the CLT."} That is mathematically false and loses credit. Proportions are justified normal by \textbf{Large Counts ($np \ge 10, n(1-p) \ge 10$)}.
  \item Furthermore, the CLT says the \textbf{sampling distribution} becomes normal as $n$ grows; it does NOT say the population or the raw sample becomes normal!
\end{itemize}
\end{trapbox}

\section{Worked Step-by-Step Examples}

\subsection*{Example 5.1: Sampling Distribution of Proportions in Quality Control}
\textbf{Problem:} A semiconductor foundry produces microchips with a defect rate of $p = 0.08$. An inspector tests a random sample of $n = 250$ microchips from a production run of 10,000 chips.
(a) Confirm that the sampling distribution of $\hat{p}$ is approximately normal.\\
(b) Calculate the mean and standard deviation of $\hat{p}$.\\
(c) What is the probability that the sample contains more than 10\% defective chips?

\textbf{Solution:}
\begin{itemize}[leftmargin=1.2em]
  \item \textbf{Part (a):}
  \textit{Random:} Stated random sample.\\
  \textit{10\% Condition:} $n = 250 \le 0.10(10,000) = 1,000$. Met.\\
  \textit{Large Counts:} $np = 250(0.08) = 20 \ge 10$ and $n(1-p) = 250(0.92) = 230 \ge 10$. Both are at least 10, so the sampling distribution of $\hat{p}$ is approximately normal.
  \item \textbf{Part (b):} $\mu_{\hat{p}} = p = \mathbf{0.08}$.\\
  $\sigma_{\hat{p}} = \sqrt{\frac{0.08(0.92)}{250}} = \sqrt{\frac{0.0736}{250}} = \sqrt{0.0002944} \approx \mathbf{0.01716}$.
  \item \textbf{Part (c):} We seek $P(\hat{p} > 0.10)$. Standardize $\hat{p} = 0.10$:
  \[ z = \frac{0.10 - 0.08}{0.01716} = \frac{0.02}{0.01716} \approx 1.17 \]
  $P(\hat{p} > 0.10) = P(Z > 1.17) = 1 - P(Z \le 1.17) = 1 - 0.8790 = \mathbf{0.1210}$ (approx. \textbf{12.1\%}).
\end{itemize}

\subsection*{Example 5.2: Applying the Central Limit Theorem to Heavily Skewed Data}
\textbf{Problem:} The claim processing times at an insurance company are heavily right-skewed with a mean of $\mu = 14.5\text{ days}$ and a standard deviation of $\sigma = 6.0\text{ days}$. A manager reviews an SRS of $n = 36$ claims.
What is the probability that the average processing time $\bar{x}$ for these 36 claims is under $13.0\text{ days}$?

\textbf{Solution:}
\begin{itemize}[leftmargin=1.2em]
  \item \textbf{Condition Check:} Although the population distribution is heavily skewed, the sample size is $n = 36 \ge 30$. By the Central Limit Theorem, the sampling distribution of the sample mean $\bar{x}$ is approximately normal!
  \item \textbf{Parameters of $\bar{x}$:}
  $\mu_{\bar{x}} = \mu = 14.5\text{ days}$.\\
  $\sigma_{\bar{x}} = \sigma / \sqrt{n} = 6.0 / \sqrt{36} = 6.0 / 6 = 1.0\text{ day}$.
  \item \textbf{Probability Calculation:} $P(\bar{x} < 13.0)$:
  \[ z = \frac{13.0 - 14.5}{1.0} = \frac{-1.5}{1.0} = -1.50 \]
  $P(\bar{x} < 13.0) = P(Z < -1.50) = \mathbf{0.0668}$ (approx. \textbf{6.68\%}).
\end{itemize}

\section{TI-84 Plus CE Calculator Playbook}

\begin{calculatorbox}[TI-84 Playbook: Sampling Distribution Probabilities]
\begin{itemize}[leftmargin=1.2em]
  \item \textbf{For Proportions $\hat{p}$:} Press \texttt{[2nd] [VARS] $\rightarrow$ 2:normalcdf(}.
  Enter \texttt{lower}, \texttt{upper}, $\mu = p$, $\sigma = \sqrt{p(1-p)/n}$.
  \item \textbf{For Means $\bar{x}$:} Press \texttt{[2nd] [VARS] $\rightarrow$ 2:normalcdf(}.
  Enter \texttt{lower}, \texttt{upper}, $\mu = \mu_{\text{pop}}$, $\sigma = \sigma_{\text{pop}}/\sqrt{n}$.
\end{itemize}
\end{calculatorbox}

\section{Comprehensive Practice Problem Set}

\subsection*{Part I: 20 Multiple-Choice Questions}

\begin{enumerate}[leftmargin=1.5em]
  \item Which of the following is a statistic?
  \begin{enumerate}[label=(\Alph*)]
    \item The true proportion $p$ of all voters \item The population mean $\mu$ \item The sample proportion $\hat{p}$ \item The population standard deviation $\sigma$ \item The census parameter
  \end{enumerate}
  \textbf{Answer: (C).} $\hat{p}$ is computed from sample data.

  \item As sample size $n$ increases from 100 to 400, the standard deviation of $\hat{p}$:
  \begin{enumerate}[label=(\Alph*)]
    \item Multiplies by 4 \item Multiplies by 2 \item Is cut in half \item Is divided by 4 \item Remains unchanged
  \end{enumerate}
  \textbf{Answer: (C).} $\sigma_{\hat{p}} \propto 1/\sqrt{n}$. $\sqrt{400}/\sqrt{100} = 20/10 = 2 \implies$ divided by 2.

  \item A population distribution is heavily skewed to the left with $\mu = 80, \sigma = 12$. A random sample of size $n = 64$ is selected. The sampling distribution of $\bar{x}$ is:
  \begin{enumerate}[label=(\Alph*)]
    \item Heavily skewed left with $\mu = 80, \sigma = 12$
    \item Approximately normal with $\mu = 80, \sigma = 1.5$
    \item Approximately normal with $\mu = 80, \sigma = 12$
    \item Uniform with $\mu = 80, \sigma = 1.5$
    \item Cannot be determined because the population is skewed
  \end{enumerate}
  \textbf{Answer: (B).} By the CLT ($n = 64 \ge 30$), the sampling distribution is approximately normal. $\sigma_{\bar{x}} = 12/\sqrt{64} = 1.5$.

  \item The Central Limit Theorem states that:
  \begin{enumerate}[label=(\Alph*)]
    \item As $n$ grows, the sample data looks more normal.
    \item As $n$ grows, the population becomes normal.
    \item As $n$ grows, the sampling distribution of $\bar{x}$ approaches a normal distribution.
    \item As $n$ grows, the standard error increases.
    \item Proportions are always normal.
  \end{enumerate}
  \textbf{Answer: (C).} Standard definition of the CLT.

  \item Which condition justifies using $\sigma_{\hat{p}} = \sqrt{p(1-p)/n}$?
  \begin{enumerate}[label=(\Alph*)]
    \item $np \ge 10$ and $n(1-p) \ge 10$ \item $n \ge 30$ \item $n \le 0.10N$ (10\% Condition) \item Population is normal \item $p = 0.50$
  \end{enumerate}
  \textbf{Answer: (C).} The 10\% condition satisfies the independence requirement for standard error formulas.

  \item An unbiased estimator is a statistic that:
  \begin{enumerate}[label=(\Alph*)]
    \item Has a standard deviation of zero.
    \item Has a sampling distribution whose mean equals the true parameter.
    \item Never produces an error.
    \item Is calculated from a census.
    \item Has a sample size greater than 100.
  \end{enumerate}
  \textbf{Answer: (B).} Defining condition of unbiasedness.

  \item If $p = 0.60$ and $n = 100$, what is $P(\hat{p} \le 0.50)$?
  \begin{enumerate}[label=(\Alph*)]
    \item 0.0207 \item 0.0500 \item 0.1000 \item 0.4793 \item 0.9793
  \end{enumerate}
  \textbf{Answer: (A).} $\sigma = \sqrt{0.6(0.4)/100} = 0.04899$. $z = (0.50 - 0.60)/0.04899 = -2.04$. $P(Z < -2.04) = 0.0207$.

  \item A population has $\mu = 50, \sigma = 10$. For an SRS of size $n = 25$, what is $P(\bar{x} > 54)$ if the population is normal?
  \begin{enumerate}[label=(\Alph*)]
    \item 0.0228 \item 0.0500 \item 0.1587 \item 0.3446 \item 0.4500
  \end{enumerate}
  \textbf{Answer: (A).} $\sigma_{\bar{x}} = 10/\sqrt{25} = 2.0$. $z = (54 - 50)/2 = 2.00$. $P(Z > 2.00) = 0.0228$.

  \item Which of the following estimators of a population mean $\mu$ is biased?
  \begin{enumerate}[label=(\Alph*)]
    \item Sample mean $\bar{x}$ \item Sample median (for symmetric populations) \item $\bar{x} + 2$ \item Midrange (for symmetric populations) \item None of the above
  \end{enumerate}
  \textbf{Answer: (C).} Adding a constant of 2 guarantees the expected value is $\mu + 2 \neq \mu$.

  \item For a population with $p = 0.02$, what is the minimum sample size required to satisfy the Large Counts condition?
  \begin{enumerate}[label=(\Alph*)]
    \item 30 \item 100 \item 250 \item 500 \item 1000
  \end{enumerate}
  \textbf{Answer: (D).} $np \ge 10 \implies n(0.02) \ge 10 \implies n \ge 10 / 0.02 = 500$.
\end{enumerate}

\begin{enumerate}[resume, leftmargin=1.5em]
  \item Why does the standard deviation of $\bar{x}$ decrease as $n$ increases?
  \begin{enumerate}[label=(\Alph*)]
    \item Larger samples average out extreme individual values, clustering sample means closer to $\mu$.
    \item The population standard deviation shrinks.
    \item The degrees of freedom decrease.
    \item The sampling frame becomes biased.
    \item $\mu$ increases.
  \end{enumerate}
  \textbf{Answer: (A).} Averaging over larger sample sizes dampens extreme individual deviations.

  \item A fair coin is tossed $n$ times. What is the standard deviation of the proportion of heads if $n = 400$?
  \begin{enumerate}[label=(\Alph*)]
    \item 0.05 \item 0.025 \item 0.0125 \item 0.25 \item 0.0025
  \end{enumerate}
  \textbf{Answer: (B).} $\sigma = \sqrt{0.5(0.5)/400} = \sqrt{0.25/400} = 0.5/20 = 0.025$.

  \item If the shape of a population is uniform between 0 and 10, the sampling distribution of $\bar{x}$ for $n = 50$ will be:
  \begin{enumerate}[label=(\Alph*)]
    \item Uniform \item Skewed right \item Approximately normal \item Bimodal \item Exponential
  \end{enumerate}
  \textbf{Answer: (C).} By the Central Limit Theorem ($n = 50 \ge 30$).

  \item The spread of a sampling distribution is also commonly called the:
  \begin{enumerate}[label=(\Alph*)]
    \item Parameter \item Standard error of the statistic \item Bias \item Confidence level \item P-value
  \end{enumerate}
  \textbf{Answer: (B).} The standard deviation of an estimator is the standard error.

  \item If $n = 16$ from a normal population with $\sigma = 8$, the standard deviation of $\bar{x}$ is:
  \begin{enumerate}[label=(\Alph*)]
    \item 8 \item 4 \item 2 \item 0.5 \item 1
  \end{enumerate}
  \textbf{Answer: (C).} $\sigma_{\bar{x}} = 8/\sqrt{16} = 8/4 = 2$.

  \item A study surveys 100 students from a university of 20,000. Is the 10\% condition satisfied?
  \begin{enumerate}[label=(\Alph*)]
    \item Yes, because $100 \le 0.10(20,000) = 2,000$.
    \item No, because $100 < 30$.
    \item Yes, because $100 \ge 30$.
    \item No, because the university is too large.
    \item Yes, by the Central Limit Theorem.
  \end{enumerate}
  \textbf{Answer: (A).} $100 \le 2,000$.

  \item What happens to the mean of a sampling distribution as the sample size quadruples?
  \begin{enumerate}[label=(\Alph*)]
    \item Quadruples \item Doubles \item Is cut in half \item Remains unchanged \item Approaches zero
  \end{enumerate}
  \textbf{Answer: (D).} The mean $\mu_{\bar{x}} = \mu$ is invariant to sample size.

  \item An estimator that consistently overestimates the true parameter is said to have:
  \begin{enumerate}[label=(\Alph*)]
    \item High variability \item Positive bias \item Low degrees of freedom \item Negative variance \item Normality
  \end{enumerate}
  \textbf{Answer: (B).} Systematically exceeding the parameter reflects positive bias.

  \item If $P(\bar{x} < 100) = 0.50$, what is the population mean $\mu$?
  \begin{enumerate}[label=(\Alph*)]
    \item 0 \item 50 \item 100 \item 200 \item Cannot be determined
  \end{enumerate}
  \textbf{Answer: (C).} Because normal distributions are symmetric around their mean, $50\%$ falls below $\mu \implies \mu = 100$.

  \item Which of the following conditions is required to apply the Central Limit Theorem?
  \begin{enumerate}[label=(\Alph*)]
    \item The population must be normal.
    \item The sample size must be sufficiently large ($n \ge 30$).
    \item $np \ge 10$.
    \item The data must be categorical.
    \item The variance must be 1.
  \end{enumerate}
  \textbf{Answer: (B).} Core criterion for the CLT.
\end{enumerate}

\subsection*{Part II: Free-Response Practice Problem}

\begin{workbookbox}[FRQ 5.1: Central Limit Theorem & Sampling Distribution of Means]
\textbf{Problem:} An automated coffee machine dispenses coffee with a mean volume of $\mu = 10.0\text{ oz}$ and a standard deviation of $\sigma = 0.8\text{ oz}$. The distribution of volumes for individual cups is skewed to the left.
(a) Can we calculate the probability that a single randomly selected cup has fewer than $9.8\text{ oz}$ of coffee? Explain.\\
(b) A tray holds an SRS of $n = 36$ cups. Describe the sampling distribution of the sample mean volume $\bar{x}$ (shape, center, spread).\\
(c) What is the probability that the average volume of these 36 cups is less than $9.8\text{ oz}$?

\textbf{Grader-Approved Score-4 Solution:}
\begin{itemize}[leftmargin=1.2em]
  \item \textbf{Part (a):} No. Because the population distribution of individual cup volumes is skewed to the left and $n=1$, we cannot assume normality for a single cup, making normal probability calculations invalid.
  \item \textbf{Part (b):}
  \textit{Shape:} Approximately Normal because $n = 36 \ge 30$, satisfying the Central Limit Theorem.\\
  \textit{Center:} $\mu_{\bar{x}} = \mu = \mathbf{10.0\text{ oz}}$.\\
  \textit{Spread:} Because $n=36 \le 0.10(\text{all cups produced})$, $\sigma_{\bar{x}} = \frac{\sigma}{\sqrt{n}} = \frac{0.8}{\sqrt{36}} = \frac{0.8}{6} \approx \mathbf{0.1333\text{ oz}}$.
  \item \textbf{Part (c):} Standardize $\bar{x} = 9.8$:
  \[ z = \frac{9.8 - 10.0}{0.1333} = \frac{-0.2}{0.1333} = -1.50 \]
  $P(\bar{x} < 9.8) = P(Z < -1.50) = \mathbf{0.0668}$ (approx. \textbf{6.68\%}).
\end{itemize}
\end{workbookbox}


\begin{frqrubricbox}[FRQ 5.2: Sampling Distribution of Means with Central Limit Theorem]
\textbf{Problem:} A commercial cargo airline transports standard shipping crates. The crate weights have a distribution with mean $\mu = 125\text{ lbs}$ and standard deviation $\sigma = 40\text{ lbs}$. The distribution of individual crate weights is heavily skewed to the right with occasional very heavy crates.
\begin{enumerate}[label=(\alph*)]
  \item Can you calculate the probability that a single randomly chosen crate weighs more than 135 lbs? Explain.
  \item A cargo aircraft is loaded with an SRS of $n = 64$ crates. Describe the sampling distribution of the sample mean crate weight $\bar{x}$, including shape, center, and spread.
  \item The maximum allowable average weight per crate for safe takeoff is 135 lbs. What is the probability that the sample mean weight of the 64 crates exceeds 135 lbs?
  \item If the airline increased the sample size to $n = 100$ crates, how would the probability of exceeding the 135-lb safety threshold change? Explain.
\end{enumerate}

\textbf{Grader-Approved Score-4 Solution \& Rubric:}
\begin{itemize}[leftmargin=1.2em]
  \item \textbf{Part (a):} No. We are told the population of individual crate weights is heavily skewed to the right. Because the population is not normal and $n = 1 < 30$, we cannot use the normal distribution to compute probabilities for a single crate.
  \item \textbf{Part (b):}
  \textit{Shape:} Approximately normal by the Central Limit Theorem because the sample size $n = 64 \ge 30$.\\
  \textit{Center:} $\mu_{\bar{x}} = \mu = 125\text{ lbs}$.\\
  \textit{Spread:} $\sigma_{\bar{x}} = \frac{\sigma}{\sqrt{n}} = \frac{40}{\sqrt{64}} = \frac{40}{8} = 5.0\text{ lbs}$ (assuming 64 crates is less than 10\% of all crates).
  \item \textbf{Part (c):} $z = \frac{135 - 125}{5.0} = \frac{10}{5.0} = 2.00$.\\
  $P(\bar{x} > 135) = P(Z > 2.00) = 1 - 0.9772 = \mathbf{0.0228}$ (approx. 2.28\%).
  \item \textbf{Part (d):} For $n = 100$, $\sigma_{\bar{x}} = \frac{40}{\sqrt{100}} = 4.0\text{ lbs}$.
  $z = \frac{135 - 125}{4.0} = 2.50$, so $P(Z > 2.50) = 0.0062$.
  Increasing $n$ reduces the standard error of the sample mean, concentrating the sampling distribution closer to $\mu = 125$ and decreasing the probability of exceeding 135 lbs.
\end{itemize}
\end{frqrubricbox}

\begin{trapbox}[Unit 5 Top 5 AP Exam Pitfalls to Avoid]
\begin{enumerate}[leftmargin=1.2em]
  \item \textbf{Claiming the CLT Makes the Population Normal:} The Central Limit Theorem only normalizes the \textbf{sampling distribution of $\bar{x}$}, NEVER the population or sample!
  \item \textbf{Applying CLT to Proportions:} CLT applies to sample means $\bar{x}$. For proportions $\hat{p}$, normality is justified by the Large Counts condition ($np \ge 10$ and $n(1-p) \ge 10$).
  \item \textbf{Forgetting to Divide by $\sqrt{n}$:} The standard deviation of $\bar{x}$ is $\sigma / \sqrt{n}$, not $\sigma$!
  \item \textbf{Confusing the 10\% Condition with Large Counts:} 10\% condition ($n \le 0.10N$) checks independence; Large Counts checks normality.
  \item \textbf{Confusing Parameter with Statistic:} $\mu$ and $p$ are fixed population parameters; $\bar{x}$ and $\hat{p}$ are variable statistics.
\end{enumerate}
\end{trapbox}
\section{Unit 5 Master Summary}
\begin{lessonbox}[Unit 5 Essential Review Checklist]
\begin{itemize}[leftmargin=1.2em]
  \item \textbf{Proportions ($\hat{p}$):} $\mu = p$, $\sigma = \sqrt{p(1-p)/n}$. Normal if $np \ge 10$ and $n(1-p) \ge 10$.
  \item \textbf{Means ($\bar{x}$):} $\mu = \mu$, $\sigma = \sigma/\sqrt{n}$. Normal if pop is normal OR $n \ge 30$ (CLT).
  \item \textbf{10\% Condition ($n \le 0.10N$):} Required for standard deviation formulas when sampling without replacement.
  \item \textbf{CLT Distinction:} CLT applies ONLY to means ($\bar{x}$), never to proportions!
\end{itemize}
\end{lessonbox}
\clearpage
